% =====================================================
% core2extreme Cybersecurity Blue Team Writeup Template
% Overleaf-ready LaTeX Template
% Author: Michael Fajkus
% Repository: https://github.com/core2extreme/writeup_guides_resources
% =====================================================

\documentclass[11pt,a4paper]{article}

% =====================================================
% PAGE SETUP
% =====================================================
\usepackage[a4paper,margin=0.9in]{geometry}
\usepackage[T1]{fontenc}
\usepackage[utf8]{inputenc}
\usepackage{graphicx}
\usepackage{xcolor}
\usepackage{titlesec}
\usepackage{hyperref}
\usepackage[most]{tcolorbox}
\usepackage{listings}
\usepackage{enumitem}
\usepackage{parskip}
\usepackage{pagecolor}
\usepackage{fancyhdr}
\usepackage{float}
\usepackage{longtable}
\usepackage{array}
\usepackage{tocloft}
\usepackage{needspace}
\usepackage{caption}
\usepackage{booktabs}
\usepackage{tabularx}

% =====================================================
% FONT SETUP
% =====================================================
\usepackage[sfdefault]{FiraSans}
\renewcommand{\familydefault}{\sfdefault}
\usepackage[scaled=0.9]{FiraMono}

% =====================================================
% TEMPLATE VARIABLES
% Change these values for every new writeup.
% =====================================================
\newcommand{\doctitle}{Cybersecurity Blue Team Writeup}
\newcommand{\docsubtitle}{Technical Guide / Short Analyst Playbook}
\newcommand{\doccategory}{SOC / DFIR / Detection Engineering}
\newcommand{\docauthor}{Your Name}
\newcommand{\docversion}{1.0}
\newcommand{\docrepo}{https://github.com/core2extreme/writeup\_guides\_resources}
\newcommand{\docrepourl}{https://github.com/core2extreme/writeup_guides_resources}
\newcommand{\doctags}{Blue Team | SOC | DFIR | Detection Engineering | Threat Hunting}

% =====================================================
% COLOR PALETTE
% =====================================================
\definecolor{bg}{HTML}{0B1020}
\definecolor{text}{HTML}{C9D1D9}
\definecolor{muted}{HTML}{8B949E}
\definecolor{primary}{HTML}{00D1FF}
\definecolor{secondary}{HTML}{7C4DFF}
\definecolor{accent}{HTML}{00FFB3}
\definecolor{warning}{HTML}{FF5F56}
\definecolor{boxbg}{HTML}{111827}
\definecolor{codebg}{HTML}{161B22}
\definecolor{ioc}{HTML}{FFD166}
\definecolor{mitre}{HTML}{FF8C42}
\definecolor{detection}{HTML}{00E5FF}
\definecolor{ir}{HTML}{B388FF}

\pagecolor{bg}
\color{text}

% =====================================================
% GLOBAL SPACING
% =====================================================
\setlength{\parskip}{0.45em}
\setlength{\parindent}{0pt}
\setlength{\textfloatsep}{8pt}
\setlength{\floatsep}{8pt}
\setlength{\intextsep}{8pt}

% =====================================================
% HYPERLINKS AND PDF METADATA
% =====================================================
\hypersetup{
    pdftitle={\doctitle},
    pdfauthor={\docauthor},
    pdfsubject={\doccategory},
    pdfkeywords={SOC, DFIR, Blue Team, Detection Engineering, Threat Hunting, Incident Response},
    colorlinks=true,
    urlcolor=primary,
    linkcolor=primary,
    citecolor=accent
}

% =====================================================
% TABLE OF CONTENTS STYLE
% =====================================================
\renewcommand{\cftsecfont}{\color{primary}\bfseries}
\renewcommand{\cftsubsecfont}{\color{text}}
\renewcommand{\cftsubsubsecfont}{\color{muted}}
\renewcommand{\cftsecpagefont}{\color{primary}\bfseries}
\renewcommand{\cftsubsecpagefont}{\color{text}}
\renewcommand{\cftsubsubsecpagefont}{\color{muted}}
\renewcommand{\cftdotsep}{1}
% =====================================================
% core2extreme Cybersecurity Blue Team Writeup Template
% Overleaf-ready LaTeX Template
% Author: Michael Fajkus
% Repository: https://github.com/core2extreme/writeup_guides_resources
% =====================================================

\documentclass[11pt,a4paper]{article}

% =====================================================
% PAGE SETUP
% =====================================================
\usepackage[a4paper,margin=0.9in]{geometry}
\usepackage[T1]{fontenc}
\usepackage[utf8]{inputenc}
\usepackage{graphicx}
\usepackage{xcolor}
\usepackage{titlesec}
\usepackage{hyperref}
\usepackage[most]{tcolorbox}
\usepackage{listings}
\usepackage{enumitem}
\usepackage{parskip}
\usepackage{pagecolor}
\usepackage{fancyhdr}
\usepackage{float}
\usepackage{longtable}
\usepackage{array}
\usepackage{tocloft}
\usepackage{needspace}
\usepackage{caption}
\usepackage{booktabs}
\usepackage{tabularx}

% =====================================================
% FONT SETUP
% =====================================================
\usepackage[sfdefault]{FiraSans}
\renewcommand{\familydefault}{\sfdefault}
\usepackage[scaled=0.9]{FiraMono}

% =====================================================
% TEMPLATE VARIABLES
% Change these values for every new writeup.
% =====================================================
\newcommand{\doctitle}{Cybersecurity Blue Team Writeup}
\newcommand{\docsubtitle}{Technical Guide / Short Analyst Playbook}
\newcommand{\doccategory}{SOC / DFIR / Detection Engineering}
\newcommand{\docauthor}{Your Name}
\newcommand{\docversion}{1.0}
\newcommand{\docrepo}{https://github.com/core2extreme/writeup\_guides\_resources}
\newcommand{\docrepourl}{https://github.com/core2extreme/writeup_guides_resources}
\newcommand{\doctags}{Blue Team | SOC | DFIR | Detection Engineering | Threat Hunting}

% =====================================================
% COLOR PALETTE
% =====================================================
\definecolor{bg}{HTML}{0B1020}
\definecolor{text}{HTML}{C9D1D9}
\definecolor{muted}{HTML}{8B949E}
\definecolor{primary}{HTML}{00D1FF}
\definecolor{secondary}{HTML}{7C4DFF}
\definecolor{accent}{HTML}{00FFB3}
\definecolor{warning}{HTML}{FF5F56}
\definecolor{boxbg}{HTML}{111827}
\definecolor{codebg}{HTML}{161B22}
\definecolor{ioc}{HTML}{FFD166}
\definecolor{mitre}{HTML}{FF8C42}
\definecolor{detection}{HTML}{00E5FF}
\definecolor{ir}{HTML}{B388FF}

\pagecolor{bg}
\color{text}

% =====================================================
% GLOBAL SPACING
% =====================================================
\setlength{\parskip}{0.45em}
\setlength{\parindent}{0pt}
\setlength{\textfloatsep}{8pt}
\setlength{\floatsep}{8pt}
\setlength{\intextsep}{8pt}

% =====================================================
% HYPERLINKS AND PDF METADATA
% =====================================================
\hypersetup{
    pdftitle={\doctitle},
    pdfauthor={\docauthor},
    pdfsubject={\doccategory},
    pdfkeywords={SOC, DFIR, Blue Team, Detection Engineering, Threat Hunting, Incident Response},
    colorlinks=true,
    urlcolor=primary,
    linkcolor=primary,
    citecolor=accent
}

% =====================================================
% TABLE OF CONTENTS STYLE
% =====================================================
\renewcommand{\cftsecfont}{\color{primary}\bfseries}
\renewcommand{\cftsubsecfont}{\color{text}}
\renewcommand{\cftsubsubsecfont}{\color{muted}}
\renewcommand{\cftsecpagefont}{\color{primary}\bfseries}
\renewcommand{\cftsubsecpagefont}{\color{text}}
\renewcommand{\cftsubsubsecpagefont}{\color{muted}}
\renewcommand{\cftdotsep}{1}

% =====================================================
% CAPTION STYLE
% =====================================================
\captionsetup{
    labelfont={bf,color=primary},
    textfont={color=text},
    font=small,
    skip=6pt
}

% =====================================================
% HEADER / FOOTER
% =====================================================
\pagestyle{fancy}
\fancyhf{}
\fancyhead[L]{\textcolor{primary}{\doctitle}}
\fancyhead[R]{\textcolor{secondary}{\doccategory}}
\fancyfoot[L]{\textcolor{muted}{core2extreme}}
\fancyfoot[C]{\textcolor{primary}{\thepage}}
\fancyfoot[R]{\textcolor{muted}{\docversion}}

\renewcommand{\headrulewidth}{0.4pt}
\renewcommand{\footrulewidth}{0.4pt}
\renewcommand{\headrule}{\hbox to\headwidth{\color{primary}\leaders\hrule height \headrulewidth\hfill}}
\renewcommand{\footrule}{\hbox to\headwidth{\color{secondary}\leaders\hrule height \footrulewidth\hfill}}

% =====================================================
% SECTION STYLE
% =====================================================
\titleformat{\section}
{\Huge\bfseries\color{primary}}
{\thesection}
{0.6em}
{}

\titleformat{\subsection}
{\Large\bfseries\color{accent}}
{\thesubsection}
{0.6em}
{}

\titleformat{\subsubsection}
{\large\bfseries\color{secondary}}
{\thesubsubsection}
{0.6em}
{}

\titlespacing*{\section}{0pt}{0.6cm}{0.45cm}
\titlespacing*{\subsection}{0pt}{0.55cm}{0.25cm}
\titlespacing*{\subsubsection}{0pt}{0.4cm}{0.2cm}

% Start every new main heading on a new page.
\let\oldsection\section
\renewcommand{\section}{\clearpage\oldsection}

% =====================================================
% LIST STYLE
% =====================================================
\setlist[itemize]{label=\textcolor{primary}{>}, leftmargin=1.5em}
\setlist[enumerate]{label=\textcolor{primary}{\arabic*.}, leftmargin=1.7em}

% =====================================================
% CODE LISTING STYLE
% =====================================================
\lstdefinestyle{cyber}{
    backgroundcolor=\color{codebg},
    basicstyle=\ttfamily\small\color{text},
    keywordstyle=\color{primary}\bfseries,
    commentstyle=\color{accent},
    stringstyle=\color{warning},
    numberstyle=\tiny\color{muted},
    breaklines=true,
    frame=single,
    rulecolor=\color{primary},
    columns=fullflexible,
    keepspaces=true,
    showstringspaces=false,
    numbers=left,
    xleftmargin=1.4em,
    framexleftmargin=1.2em
}

\lstset{style=cyber}

% =====================================================
% CUSTOM BOXES
% =====================================================
\newtcolorbox{socbox}[1][]{
    enhanced,
    colback=boxbg,
    colframe=primary,
    coltext=text,
    arc=6pt,
    boxrule=1pt,
    left=10pt,
    right=10pt,
    top=10pt,
    bottom=10pt,
    title=\textbf{SOC Note},
    fonttitle=\bfseries\color{bg},
    coltitle=bg,
    colbacktitle=primary,
    #1
}

\newtcolorbox{warningbox}[1][]{
    enhanced,
    colback=boxbg,
    colframe=warning,
    coltext=text,
    arc=6pt,
    boxrule=1pt,
    left=10pt,
    right=10pt,
    top=10pt,
    bottom=10pt,
    title=\textbf{Warning},
    fonttitle=\bfseries\color{bg},
    coltitle=bg,
    colbacktitle=warning,
    #1
}

\newtcolorbox{detectionbox}[1][]{
    enhanced,
    colback=boxbg,
    colframe=detection,
    coltext=text,
    arc=6pt,
    boxrule=1pt,
    left=10pt,
    right=10pt,
    top=10pt,
    bottom=10pt,
    title=\textbf{Detection Opportunity},
    fonttitle=\bfseries\color{bg},
    coltitle=bg,
    colbacktitle=detection,
    #1
}

\newtcolorbox{iocbox}[1][]{
    enhanced,
    colback=boxbg,
    colframe=ioc,
    coltext=text,
    arc=6pt,
    boxrule=1pt,
    left=10pt,
    right=10pt,
    top=10pt,
    bottom=10pt,
    title=\textbf{IOC Box},
    fonttitle=\bfseries\color{bg},
    coltitle=bg,
    colbacktitle=ioc,
    #1
}

\newtcolorbox{mitrebox}[1][]{
    enhanced,
    colback=boxbg,
    colframe=mitre,
    coltext=text,
    arc=6pt,
    boxrule=1pt,
    left=10pt,
    right=10pt,
    top=10pt,
    bottom=10pt,
    title=\textbf{MITRE ATT\&CK Mapping},
    fonttitle=\bfseries\color{bg},
    coltitle=bg,
    colbacktitle=mitre,
    #1
}

\newtcolorbox{irbox}[1][]{
    enhanced,
    colback=boxbg,
    colframe=ir,
    coltext=text,
    arc=6pt,
    boxrule=1pt,
    left=10pt,
    right=10pt,
    top=10pt,
    bottom=10pt,
    title=\textbf{IR Tip},
    fonttitle=\bfseries\color{bg},
    coltitle=bg,
    colbacktitle=ir,
    #1
}

% =====================================================
% HELPER COMMANDS
% =====================================================
\newcommand{\placeholder}{\textcolor{muted}{Lorem ipsum dolor sit amet, consectetur adipiscing elit. Replace this placeholder with writeup content.}}
\newcommand{\inlinecode}[1]{\texttt{\textcolor{accent}{#1}}}

% Generic figure placeholder.
% Usage:
% \imageplaceholder{Caption text}
\newcommand{\imageplaceholder}[1]{
\begin{figure}[H]
\centering
\begin{tcolorbox}[
    colback=codebg,
    colframe=primary,
    coltext=muted,
    width=\textwidth,
    height=0.28\textheight,
    valign=center,
    halign=center,
    arc=6pt,
    boxrule=1pt
]
Image / Diagram Placeholder
\end{tcolorbox}
\caption{#1}
\end{figure}
}

% =====================================================
% TITLE PAGE
% =====================================================
\title{
\vspace{-1cm}

\Huge
\textbf{\textcolor{primary}{Title 1}}\\

\vspace{0.2cm}

\textbf{\textcolor{white}{\doctitle}}

\vspace{0.8cm}

\Large
\textcolor{accent}{\docsubtitle}

\vspace{0.8cm}

\normalsize
\textcolor{muted}{\doctags}

\vspace{1.2cm}

\Large
\textcolor{text}{\docauthor}

\vspace{0.3cm}

\normalsize
\textcolor{muted}{Version \docversion{} --- \today}
}

\author{}
\date{}

% =====================================================
% DOCUMENT START
% =====================================================
\begin{document}

\pagenumbering{gobble}

\maketitle
\thispagestyle{empty}

\vfill

\begin{center}
\begin{tcolorbox}[
    colback=boxbg,
    colframe=primary,
    coltext=text,
    width=0.85\textwidth,
    arc=6pt,
    boxrule=1pt,
    left=12pt,
    right=12pt,
    top=12pt,
    bottom=12pt
]
\centering
\textbf{\textcolor{primary}{Purpose}}\\
Reusable Blue Team writeup template for SOC analysts, DFIR practitioners, incident responders, threat hunters and detection engineers.
\end{tcolorbox}
\end{center}

\clearpage

\tableofcontents
\thispagestyle{empty}

\clearpage

\pagenumbering{arabic}
\setcounter{page}{1}

% =====================================================
% CONTENT SKELETON
% =====================================================

\section{Heading 1}

\begin{socbox}
\placeholder
\end{socbox}

\subsection{Subheading 1.1}

\placeholder

\subsection{Subheading 1.2}

\placeholder

\begin{lstlisting}[language=bash,caption={Command Example}]
# Replace this with a real command, query, rule, or script.
echo "Hello Blue Team"
\end{lstlisting}

\section{Heading 2}

\subsection{Subheading 2.1}

\placeholder

\subsection{Subheading 2.2}

\placeholder

\begin{detectionbox}
Use this box for detection logic, hunting ideas, SIEM queries, analytic notes or suspicious behavior patterns.
\end{detectionbox}

\section{Heading 3}

\subsection{Subheading 3.1}

\placeholder

\subsection{Subheading 3.2}

\placeholder

\imageplaceholder{Generic Diagram Placeholder}

\section{Heading 4}

\subsection{Subheading 4.1}

\placeholder

\subsection{Subheading 4.2}

\placeholder

\begin{iocbox}
\begin{itemize}
    \item Domain: \inlinecode{example.com}
    \item IP Address: \inlinecode{203.0.113.10}
    \item URL: \inlinecode{https://example.com/path}
    \item Hash: \inlinecode{aaaaaaaaaaaaaaaaaaaaaaaaaaaaaaaaaaaaaaaaaaaaaaaaaaaaaaaaaaaaaaaa}
\end{itemize}
\end{iocbox}

\section{Heading 5}

\subsection{Subheading 5.1}

\placeholder

\subsection{Subheading 5.2}

\placeholder

\begin{mitrebox}
\begin{itemize}
    \item Technique ID: \inlinecode{TXXXX}
    \item Technique Name: Placeholder Technique
    \item Tactic: Placeholder Tactic
    \item Analyst Note: Replace with real ATT\&CK context.
\end{itemize}
\end{mitrebox}

\section{Heading 6}

\subsection{Subheading 6.1}

\placeholder

\subsection{Subheading 6.2}

\placeholder

\begin{warningbox}
Use this box for caveats, false positives, operational risks, assumptions or investigation warnings.
\end{warningbox}

\section{Heading 7}

\subsection{Subheading 7.1}

\placeholder

\subsection{Subheading 7.2}

\placeholder

\begin{irbox}
Use this box for containment, eradication, recovery, escalation or evidence preservation notes.
\end{irbox}

\section{Heading 8}

\subsection{Subheading 8.1}

\placeholder

\subsection{Subheading 8.2}

\placeholder

\begin{longtable}{>{\raggedright\arraybackslash}p{0.28\textwidth} >{\raggedright\arraybackslash}p{0.62\textwidth}}
\toprule
\textcolor{primary}{\textbf{Field}} & \textcolor{primary}{\textbf{Description}} \\
\midrule
Placeholder 1 & Replace with analyst-relevant information. \\
Placeholder 2 & Replace with analyst-relevant information. \\
Placeholder 3 & Replace with analyst-relevant information. \\
\bottomrule
\end{longtable}

\section{Heading 9}

\subsection{Subheading 9.1}

\placeholder

\subsection{Subheading 9.2}

\placeholder

\begin{lstlisting}[language=Python,caption={Python Example}]
# Replace this with real analysis code.
def analyze_artifact(artifact):
    return f"Analyzing {artifact}"

print(analyze_artifact("placeholder"))
\end{lstlisting}

\section{Heading 10}

\subsection{Subheading 10.1}

\placeholder

\subsection{Subheading 10.2}

\placeholder

\begin{socbox}
Use this final note for key takeaways, summary points or analyst recommendations.
\end{socbox}

% =====================================================
% RESOURCES
% =====================================================
\section{Resources}

\begin{socbox}
\textbf{GitHub Repository:}

\href{\docrepourl}{\docrepo}

\vspace{0.3cm}

Repository contents may include:

\begin{itemize}
    \item Full-resolution diagrams
    \item PDF exports
    \item LaTeX source
    \item IOC examples
    \item Sample logs
    \item Detection examples
    \item Additional references
\end{itemize}
\end{socbox}

\subsection{Additional References}

\begin{itemize}
    \item Vendor documentation
    \item Microsoft Learn
    \item MITRE ATT\&CK
    \item Security research blogs
    \item Internal SOC / DFIR notes
\end{itemize}

% =====================================================
% DISCLAIMER
% =====================================================
\section{Disclaimer}

\begin{warningbox}
Educational and defensive security purposes only.

All content is intended for Blue Team, Detection Engineering, Threat Hunting, Incident Response and Security Operations use cases.
\end{warningbox}

\vspace{1cm}

\begin{center}
\Huge
\textcolor{primary}{Stay curious. Hunt anomalies.}
\end{center}

\end{document}

% =====================================================
% CAPTION STYLE
% =====================================================
\captionsetup{
    labelfont={bf,color=primary},
    textfont={color=text},
    font=small,
    skip=6pt
}

% =====================================================
% HEADER / FOOTER
% =====================================================
\pagestyle{fancy}
\fancyhf{}
\fancyhead[L]{\textcolor{primary}{\doctitle}}
\fancyhead[R]{\textcolor{secondary}{\doccategory}}
\fancyfoot[L]{\textcolor{muted}{core2extreme}}
\fancyfoot[C]{\textcolor{primary}{\thepage}}
\fancyfoot[R]{\textcolor{muted}{\docversion}}

\renewcommand{\headrulewidth}{0.4pt}
\renewcommand{\footrulewidth}{0.4pt}
\renewcommand{\headrule}{\hbox to\headwidth{\color{primary}\leaders\hrule height \headrulewidth\hfill}}
\renewcommand{\footrule}{\hbox to\headwidth{\color{secondary}\leaders\hrule height \footrulewidth\hfill}}

% =====================================================
% SECTION STYLE
% =====================================================
\titleformat{\section}
{\Huge\bfseries\color{primary}}
{\thesection}
{0.6em}
{}

\titleformat{\subsection}
{\Large\bfseries\color{accent}}
{\thesubsection}
{0.6em}
{}

\titleformat{\subsubsection}
{\large\bfseries\color{secondary}}
{\thesubsubsection}
{0.6em}
{}

\titlespacing*{\section}{0pt}{0.6cm}{0.45cm}
\titlespacing*{\subsection}{0pt}{0.55cm}{0.25cm}
\titlespacing*{\subsubsection}{0pt}{0.4cm}{0.2cm}

% Start every new main heading on a new page.
\let\oldsection\section
\renewcommand{\section}{\clearpage\oldsection}

% =====================================================
% LIST STYLE
% =====================================================
\setlist[itemize]{label=\textcolor{primary}{>}, leftmargin=1.5em}
\setlist[enumerate]{label=\textcolor{primary}{\arabic*.}, leftmargin=1.7em}

% =====================================================
% CODE LISTING STYLE
% =====================================================
\lstdefinestyle{cyber}{
    backgroundcolor=\color{codebg},
    basicstyle=\ttfamily\small\color{text},
    keywordstyle=\color{primary}\bfseries,
    commentstyle=\color{accent},
    stringstyle=\color{warning},
    numberstyle=\tiny\color{muted},
    breaklines=true,
    frame=single,
    rulecolor=\color{primary},
    columns=fullflexible,
    keepspaces=true,
    showstringspaces=false,
    numbers=left,
    xleftmargin=1.4em,
    framexleftmargin=1.2em
}

\lstset{style=cyber}

% =====================================================
% CUSTOM BOXES
% =====================================================
\newtcolorbox{socbox}[1][]{
    enhanced,
    colback=boxbg,
    colframe=primary,
    coltext=text,
    arc=6pt,
    boxrule=1pt,
    left=10pt,
    right=10pt,
    top=10pt,
    bottom=10pt,
    title=\textbf{SOC Note},
    fonttitle=\bfseries\color{bg},
    coltitle=bg,
    colbacktitle=primary,
    #1
}

\newtcolorbox{warningbox}[1][]{
    enhanced,
    colback=boxbg,
    colframe=warning,
    coltext=text,
    arc=6pt,
    boxrule=1pt,
    left=10pt,
    right=10pt,
    top=10pt,
    bottom=10pt,
    title=\textbf{Warning},
    fonttitle=\bfseries\color{bg},
    coltitle=bg,
    colbacktitle=warning,
    #1
}

\newtcolorbox{detectionbox}[1][]{
    enhanced,
    colback=boxbg,
    colframe=detection,
    coltext=text,
    arc=6pt,
    boxrule=1pt,
    left=10pt,
    right=10pt,
    top=10pt,
    bottom=10pt,
    title=\textbf{Detection Opportunity},
    fonttitle=\bfseries\color{bg},
    coltitle=bg,
    colbacktitle=detection,
    #1
}

\newtcolorbox{iocbox}[1][]{
    enhanced,
    colback=boxbg,
    colframe=ioc,
    coltext=text,
    arc=6pt,
    boxrule=1pt,
    left=10pt,
    right=10pt,
    top=10pt,
    bottom=10pt,
    title=\textbf{IOC Box},
    fonttitle=\bfseries\color{bg},
    coltitle=bg,
    colbacktitle=ioc,
    #1
}

\newtcolorbox{mitrebox}[1][]{
    enhanced,
    colback=boxbg,
    colframe=mitre,
    coltext=text,
    arc=6pt,
    boxrule=1pt,
    left=10pt,
    right=10pt,
    top=10pt,
    bottom=10pt,
    title=\textbf{MITRE ATT\&CK Mapping},
    fonttitle=\bfseries\color{bg},
    coltitle=bg,
    colbacktitle=mitre,
    #1
}

\newtcolorbox{irbox}[1][]{
    enhanced,
    colback=boxbg,
    colframe=ir,
    coltext=text,
    arc=6pt,
    boxrule=1pt,
    left=10pt,
    right=10pt,
    top=10pt,
    bottom=10pt,
    title=\textbf{IR Tip},
    fonttitle=\bfseries\color{bg},
    coltitle=bg,
    colbacktitle=ir,
    #1
}

% =====================================================
% HELPER COMMANDS
% =====================================================
\newcommand{\placeholder}{\textcolor{muted}{Lorem ipsum dolor sit amet, consectetur adipiscing elit. Replace this placeholder with writeup content.}}
\newcommand{\inlinecode}[1]{\texttt{\textcolor{accent}{#1}}}

% Generic figure placeholder.
% Usage:
% \imageplaceholder{Caption text}
\newcommand{\imageplaceholder}[1]{
\begin{figure}[H]
\centering
\begin{tcolorbox}[
    colback=codebg,
    colframe=primary,
    coltext=muted,
    width=\textwidth,
    height=0.28\textheight,
    valign=center,
    halign=center,
    arc=6pt,
    boxrule=1pt
]
Image / Diagram Placeholder
\end{tcolorbox}
\caption{#1}
\end{figure}
}

% =====================================================
% TITLE PAGE
% =====================================================
\title{
\vspace{-1cm}

\Huge
\textbf{\textcolor{primary}{core2extreme}}\\

\vspace{0.2cm}

\textbf{\textcolor{white}{\doctitle}}

\vspace{0.8cm}

\Large
\textcolor{accent}{\docsubtitle}

\vspace{0.8cm}

\normalsize
\textcolor{muted}{\doctags}

\vspace{1.2cm}

\Large
\textcolor{text}{\docauthor}

\vspace{0.3cm}

\normalsize
\textcolor{muted}{Version \docversion{} --- \today}
}

\author{}
\date{}

% =====================================================
% DOCUMENT START
% =====================================================
\begin{document}

\pagenumbering{gobble}

\maketitle
\thispagestyle{empty}

\vfill

\begin{center}
\begin{tcolorbox}[
    colback=boxbg,
    colframe=primary,
    coltext=text,
    width=0.85\textwidth,
    arc=6pt,
    boxrule=1pt,
    left=12pt,
    right=12pt,
    top=12pt,
    bottom=12pt
]
\centering
\textbf{\textcolor{primary}{Purpose}}\\
Reusable Blue Team writeup template for SOC analysts, DFIR practitioners, incident responders, threat hunters and detection engineers.
\end{tcolorbox}
\end{center}

\clearpage

\tableofcontents
\thispagestyle{empty}

\clearpage

\pagenumbering{arabic}
\setcounter{page}{1}

% =====================================================
% CONTENT SKELETON
% =====================================================

\section{Heading 1}

\begin{socbox}
\placeholder
\end{socbox}

\subsection{Subheading 1.1}

\placeholder

\subsection{Subheading 1.2}

\placeholder

\begin{lstlisting}[language=bash,caption={Command Example}]
# Replace this with a real command, query, rule, or script.
echo "Hello Blue Team"
\end{lstlisting}

\section{Heading 2}

\subsection{Subheading 2.1}

\placeholder

\subsection{Subheading 2.2}

\placeholder

\begin{detectionbox}
Use this box for detection logic, hunting ideas, SIEM queries, analytic notes or suspicious behavior patterns.
\end{detectionbox}

\section{Heading 3}

\subsection{Subheading 3.1}

\placeholder

\subsection{Subheading 3.2}

\placeholder

\imageplaceholder{Generic Diagram Placeholder}

\section{Heading 4}

\subsection{Subheading 4.1}

\placeholder

\subsection{Subheading 4.2}

\placeholder

\begin{iocbox}
\begin{itemize}
    \item Domain: \inlinecode{example.com}
    \item IP Address: \inlinecode{203.0.113.10}
    \item URL: \inlinecode{https://example.com/path}
    \item Hash: \inlinecode{aaaaaaaaaaaaaaaaaaaaaaaaaaaaaaaaaaaaaaaaaaaaaaaaaaaaaaaaaaaaaaaa}
\end{itemize}
\end{iocbox}

\section{Heading 5}

\subsection{Subheading 5.1}

\placeholder

\subsection{Subheading 5.2}

\placeholder

\begin{mitrebox}
\begin{itemize}
    \item Technique ID: \inlinecode{TXXXX}
    \item Technique Name: Placeholder Technique
    \item Tactic: Placeholder Tactic
    \item Analyst Note: Replace with real ATT\&CK context.
\end{itemize}
\end{mitrebox}

\section{Heading 6}

\subsection{Subheading 6.1}

\placeholder

\subsection{Subheading 6.2}

\placeholder

\begin{warningbox}
Use this box for caveats, false positives, operational risks, assumptions or investigation warnings.
\end{warningbox}

\section{Heading 7}

\subsection{Subheading 7.1}

\placeholder

\subsection{Subheading 7.2}

\placeholder

\begin{irbox}
Use this box for containment, eradication, recovery, escalation or evidence preservation notes.
\end{irbox}

\section{Heading 8}

\subsection{Subheading 8.1}

\placeholder

\subsection{Subheading 8.2}

\placeholder

\begin{longtable}{>{\raggedright\arraybackslash}p{0.28\textwidth} >{\raggedright\arraybackslash}p{0.62\textwidth}}
\toprule
\textcolor{primary}{\textbf{Field}} & \textcolor{primary}{\textbf{Description}} \\
\midrule
Placeholder 1 & Replace with analyst-relevant information. \\
Placeholder 2 & Replace with analyst-relevant information. \\
Placeholder 3 & Replace with analyst-relevant information. \\
\bottomrule
\end{longtable}

\section{Heading 9}

\subsection{Subheading 9.1}

\placeholder

\subsection{Subheading 9.2}

\placeholder

\begin{lstlisting}[language=Python,caption={Python Example}]
# Replace this with real analysis code.
def analyze_artifact(artifact):
    return f"Analyzing {artifact}"

print(analyze_artifact("placeholder"))
\end{lstlisting}

\section{Heading 10}

\subsection{Subheading 10.1}

\placeholder

\subsection{Subheading 10.2}

\placeholder

\begin{socbox}
Use this final note for key takeaways, summary points or analyst recommendations.
\end{socbox}

% =====================================================
% RESOURCES
% =====================================================
\section{Resources}

\begin{socbox}
\textbf{GitHub Repository:}

\href{\docrepourl}{\docrepo}

\vspace{0.3cm}

Repository contents may include:

\begin{itemize}
    \item Full-resolution diagrams
    \item PDF exports
    \item LaTeX source
    \item IOC examples
    \item Sample logs
    \item Detection examples
    \item Additional references
\end{itemize}
\end{socbox}

\subsection{Additional References}

\begin{itemize}
    \item Vendor documentation
    \item Microsoft Learn
    \item MITRE ATT\&CK
    \item Security research blogs
    \item Internal SOC / DFIR notes
\end{itemize}

% =====================================================
% DISCLAIMER
% =====================================================
\section{Disclaimer}

\begin{warningbox}
Educational and defensive security purposes only.

All content is intended for Blue Team, Detection Engineering, Threat Hunting, Incident Response and Security Operations use cases.
\end{warningbox}

\vspace{1cm}

\begin{center}
\Huge
\textcolor{primary}{Stay curious. Hunt anomalies.}
\end{center}

\end{document}
